% Created 2026-10-16 Fri 13:23
% Intended LaTeX compiler: pdflatex
\documentclass[11pt,letterpaper]{article}
\usepackage[utf8]{inputenc}
\usepackage[T1]{fontenc}
\usepackage[spanish]{babel}
\usepackage{graphicx}
\usepackage{amsmath}
\usepackage{amssymb}
\usepackage{hyperref}
\newcommand{\materia}{MODELADO Y PROGRAMACIÓN}
\newcommand{\tipoPractica}{DOCUMENTACIÓN}
\newcommand{\practicaTitulo}{PRÁCTICA 4:\\PRUEBAS UNITARIAS}
\newcommand{\practicaSubtitulo}{DISEÑO DE PRUEBAS PARA UN SISTEMA}
\newcommand{\autorUno}{Alan Gael Alvarez Hernández}
\newcommand{\cuentaUno}{323113340}
\newcommand{\autorDos}{María Fernanda Miranda Sánchez}
\newcommand{\cuentaDos}{420102320}
\newcommand{\autorTres}{Pedro Pablo Ruiz Hernández}
\newcommand{\cuentaTres}{323198507}
\newcommand{\nombreProfesora}{Karla Adriana Esquivel Guzmán}
\newcommand{\ayudanteTeoria}{Ossiel Aguilar Espíndola}
\newcommand{\ayudanteLaboratorio}{Enoc Pineda Aguirre}
\newcommand{\semestreFinal}{SEMESTRE 1-2027}
\newcommand{\fechaPortada}{Ciudad de México, 20 de octubre de 2026}
\newcommand{\leyendaEscudoLineaUno}{FACULTAD DE CIENCIAS}
\newcommand{\leyendaEscudoLineaDos}{UNAM}
\newcommand{\encabezadoIzquierdo}{Modelado y Programación}
\newcommand{\encabezadoDerecho}{Práctica 4: Pruebas unitarias}
\newif\ifportada
\portadatrue
\newif\ifbibliografia
\bibliografiatrue
\usepackage[T1]{fontenc}
\usepackage[utf8]{inputenc}
\usepackage{lmodern}
\usepackage{microtype}
\usepackage[normalem]{ulem}
\newcommand{\fuentePortada}{\fontfamily{ptm}\selectfont}
\newcommand{\fuenteInstitucional}{\fontfamily{ptm}\fontseries{m}\selectfont}
\AtBeginDocument{\renewcommand{\contentsname}{Índice}\renewcommand{\tablename}{Tabla}\renewcommand{\figurename}{Figura}\renewcommand{\refname}{Bibliografía}}
\AtBeginDocument{\decimalpoint}
\AtBeginDocument{\shorthandoff{"}}
\usepackage[top=2.2cm,bottom=2.2cm,left=2.3cm,right=2.3cm]{geometry}
\setlength{\parindent}{0cm}
\setlength{\parskip}{0.65em}
\emergencystretch=2em
\usepackage{graphicx}
\usepackage{float}
\usepackage{caption}
\usepackage{array}
\usepackage{longtable}
\usepackage{ragged2e}
\usepackage{makecell}
\newcolumntype{L}[1]{>{\RaggedRight\arraybackslash\sloppy\hspace{0pt}}p{#1}}
\newcolumntype{C}[1]{>{\Centering\arraybackslash}p{#1}}
\usepackage{enumitem}
\setlist[itemize,1]{label=\textbullet}
\setlist[itemize,2]{label=--}
\setlist[itemize,3]{label=$\circ$}
\setlist[enumerate,1]{label=\arabic*.}
\setlist[enumerate,2]{label=\alph*)}
\usepackage{xcolor}
\usepackage{tikz}
\usetikzlibrary{calc}
\usepackage{listings}
\lstdefinestyle{practica}{
basicstyle=\ttfamily\small,
backgroundcolor=\color{gray!7},
frame=single,
rulecolor=\color{black},
framerule=0.45pt,
framesep=5pt,
xleftmargin=0pt,
xrightmargin=0pt,
aboveskip=0.4em,
belowskip=0.8em,
showstringspaces=false,
columns=fullflexible,
keepspaces=true,
breaklines=true,
literate={á}{{\'a}}1 {é}{{\'e}}1 {í}{{\'i}}1 {ó}{{\'o}}1 {ú}{{\'u}}1 {Á}{{\'A}}1 {É}{{\'E}}1 {Í}{{\'I}}1 {Ó}{{\'O}}1 {Ú}{{\'U}}1 {ñ}{{\~n}}1 {Ñ}{{\~N}}1
}
\usepackage{titlesec}
\setcounter{secnumdepth}{2}
\setcounter{tocdepth}{2}
\titleformat{\section}{\Large\bfseries}{}{0pt}{}
\titleformat{\subsection}{\large\bfseries}{}{0pt}{}
\titlespacing*{\section}{0pt}{1.5em}{0.8em}
\titlespacing*{\subsection}{0pt}{1.2em}{0.5em}
\usepackage{fancyhdr}
\setlength{\headheight}{15pt}
\pagestyle{fancy}
\fancyhf{}
\fancyhead[L]{\encabezadoIzquierdo}
\fancyhead[R]{\encabezadoDerecho}
\fancyfoot[C]{\thepage}
\renewcommand{\headrulewidth}{0.4pt}
\renewcommand{\footrulewidth}{0pt}
\hypersetup{colorlinks=true,linkcolor=black,urlcolor=blue,citecolor=black}
\author{Nombre del alumno}
\date{\today}
\title{Reporte}
\hypersetup{
 pdfauthor={Nombre del alumno},
 pdftitle={Reporte},
 pdfkeywords={},
 pdfsubject={},
 pdfcreator={Emacs 30.1 (Org mode 9.7.11)}, 
 pdflang={Spanish}}
\begin{document}


\ifportada

\hypersetup{pageanchor=false}

\begin{titlepage}

\thispagestyle{empty}


\begin{tikzpicture}[remember picture,overlay]


% ============================================================
% ESCUDO UNAM
% ============================================================

\IfFileExists{imagenes/unam.png}{

    \node[
        anchor=north west,
        inner sep=0pt
    ]
    at ([xshift=2.92cm,yshift=-2.17cm]current page.north west)
    {
        \includegraphics[width=3.68cm]{imagenes/unam.png}
    };

}{}


% ============================================================
% UNIVERSIDAD NACIONAL AUTÓNOMA DE MÉXICO
%
% La referencia usa Times con iniciales manualmente más grandes.
% Iniciales: ~16 pt.
% Resto: ~13 pt.
% ============================================================

\node[
    anchor=north,
    align=center,
    font=\fuenteInstitucional
]
at ([xshift=12.10cm,yshift=-2.51cm]current page.north west)
{
    {\fontsize{16.02}{18.4}\selectfont U}{\fontsize{13.02}{18.4}\selectfont NIVERSIDAD }
    {\fontsize{16.02}{18.4}\selectfont N}{\fontsize{13.02}{18.4}\selectfont ACIONAL }
    {\fontsize{16.02}{18.4}\selectfont A}{\fontsize{13.02}{18.4}\selectfont UTÓNOMA}\\[0.12cm]
    {\fontsize{13.02}{18.4}\selectfont DE }
    {\fontsize{16.02}{18.4}\selectfont M}{\fontsize{13.02}{18.4}\selectfont ÉXICO}
};


% ============================================================
% DOBLE LÍNEA HORIZONTAL
% ============================================================

% Línea principal

\draw[line width=1.90pt]
([xshift=7.00cm,yshift=-3.98cm]current page.north west)
--
([xshift=17.05cm,yshift=-3.98cm]current page.north west);


% Línea secundaria

\draw[line width=0.60pt]
([xshift=7.00cm,yshift=-4.50cm]current page.north west)
--
([xshift=17.05cm,yshift=-4.50cm]current page.north west);


% ============================================================
% FACULTAD DE CIENCIAS
%
% Misma construcción tipográfica de la referencia:
% iniciales grandes + resto ligeramente menor.
% ============================================================

\node[
    anchor=north,
    align=center,
    font=\fuenteInstitucional
]
at ([xshift=12.10cm,yshift=-5.31cm]current page.north west)
{
    {\fontsize{16.02}{18.0}\selectfont F}{\fontsize{13.02}{18.0}\selectfont ACULTAD DE }
    {\fontsize{16.02}{18.0}\selectfont C}{\fontsize{13.02}{18.0}\selectfont IENCIAS}
};


% ============================================================
% TRES LÍNEAS VERTICALES
% ============================================================

% Línea izquierda

\draw[line width=0.60pt]
([xshift=4.22cm,yshift=-6.27cm]current page.north west)
--
([xshift=4.22cm,yshift=-21.35cm]current page.north west);


% Línea central

\draw[line width=1.95pt]
([xshift=4.78cm,yshift=-6.27cm]current page.north west)
--
([xshift=4.78cm,yshift=-21.35cm]current page.north west);


% Línea derecha

\draw[line width=0.60pt]
([xshift=5.35cm,yshift=-6.27cm]current page.north west)
--
([xshift=5.35cm,yshift=-21.35cm]current page.north west);


% ============================================================
% MATERIA
% ============================================================

\node[
    anchor=north,
    align=center,
    text width=11.5cm,
    font=\fuentePortada\fontsize{13.0}{15.2}\selectfont
]
at ([xshift=12.10cm,yshift=-7.78cm]current page.north west)
{
    \materia
};


% ============================================================
% TIPO DE PRÁCTICA
% ============================================================

\node[
    anchor=north,
    align=center,
    font=\fuentePortada\fontsize{10.8}{12.8}\selectfont
]
at ([xshift=12.10cm,yshift=-8.60cm]current page.north west)
{
    \tipoPractica
};


% ============================================================
% TÍTULO DE LA PRÁCTICA
% ============================================================

\node[
    anchor=north,
    align=center,
    text width=11.8cm,
    font=\fuentePortada\fontsize{20.8}{24.2}\selectfont
]
at ([xshift=12.10cm,yshift=-10.36cm]current page.north west)
{
    \practicaTitulo
};


% ============================================================
% SUBTÍTULO DE LA PRÁCTICA
% ============================================================

\node[
    anchor=north,
    align=center,
    text width=11.6cm,
    font=\fuentePortada\fontsize{10.8}{12.8}\selectfont
]
at ([xshift=12.10cm,yshift=-12.42cm]current page.north west)
{
    \practicaSubtitulo
};


% ============================================================
% PRESENTAN
% ============================================================

\node[
    anchor=north,
    align=center,
    font=\fuentePortada\fontsize{11.3}{13.4}\selectfont
]
at ([xshift=12.10cm,yshift=-14.42cm]current page.north west)
{
    P\hspace{0.31cm}
    R\hspace{0.31cm}
    E\hspace{0.31cm}
    S\hspace{0.31cm}
    E\hspace{0.31cm}
    N\hspace{0.31cm}
    T\hspace{0.31cm}
    A\hspace{0.31cm}
    N
    \hspace{0.31cm}:
};


% ============================================================
% EQUIPO
%
% Alineado con las líneas horizontales superiores:
% ancho = 10.05 cm.
%
% Separaciones:
% encabezado -> primer integrante = 0.52 cm
% entre integrantes = 0.44 cm
% último integrante -> PROFESORA = 1.04 cm
%
% Los datos se editan al inicio del archivo:
% \autorUno / \cuentaUno
% \autorDos / \cuentaDos
% \autorTres / \cuentaTres
% ============================================================

% Encabezados
% Mismo tamaño y altura tipográfica que PROFESORA:
\node[
    anchor=north,
    inner sep=0pt,
    font=\fuentePortada\mdseries\fontsize{10.8}{12.8}\selectfont
]
at ([xshift=12.025cm,yshift=-16.25cm]current page.north west)
{
    \makebox[10.05cm][s]{INTEGRANTE DEL EQUIPO:\hfill NÚM. DE CUENTA:}
};


% Integrante 1
% La separación visual EQUIPO -> primer integrante es la misma
% que PROFESORA -> nombre de profesora.
\node[
    anchor=north,
    inner sep=0pt,
    font=\fuentePortada\mdseries\fontsize{9.8}{11.6}\selectfont
]
at ([xshift=12.025cm,yshift=-16.9cm]current page.north west)
{
    \makebox[10.05cm][s]{%
        \autorUno
        \leaders\hbox{\kern0.10em.\kern0.10em}\hfill
        \cuentaUno
    }
};


% Integrante 2
\node[
    anchor=north,
    inner sep=0pt,
    font=\fuentePortada\mdseries\fontsize{9.8}{11.6}\selectfont
]
at ([xshift=12.025cm,yshift=-17.34cm]current page.north west)
{
    \makebox[10.05cm][s]{%
        \autorDos
        \leaders\hbox{\kern0.10em.\kern0.10em}\hfill
        \cuentaDos
    }
};


% Integrante 3
\node[
    anchor=north,
    inner sep=0pt,
    font=\fuentePortada\mdseries\fontsize{9.8}{11.6}\selectfont
]
at ([xshift=12.025cm,yshift=-17.78cm]current page.north west)
{
    \makebox[10.05cm][s]{%
        \autorTres
        \leaders\hbox{\kern0.10em.\kern0.10em}\hfill
        \cuentaTres
    }
};


% ============================================================
% PROFESORA
% ============================================================

\node[
    anchor=north,
    align=center,
    font=\fuentePortada\fontsize{10.8}{12.8}\selectfont
]
at ([xshift=12.10cm,yshift=-18.82cm]current page.north west)
{
    PROFESORA:
};

\node[
    anchor=north,
    align=center,
    text width=12.2cm,
    font=\fuentePortada\mdseries\fontsize{9.8}{11.6}\selectfont
]
at ([xshift=12.10cm,yshift=-19.32cm]current page.north west)
{
    \nombreProfesora
};


% ============================================================
% AYUDANTE DE TEORÍA
% ============================================================

\node[
    anchor=north,
    align=center,
    font=\fuentePortada\fontsize{10.0}{11.9}\selectfont
]
at ([xshift=12.10cm,yshift=-20.36cm]current page.north west)
{
    AYUDANTE DE TEORÍA:
};

\node[
    anchor=north,
    align=center,
    text width=12.2cm,
    font=\fuentePortada\mdseries\fontsize{9.8}{11.6}\selectfont
]
at ([xshift=12.10cm,yshift=-20.88cm]current page.north west)
{
    \ayudanteTeoria
};


% ============================================================
% AYUDANTE DE LABORATORIO
% ============================================================

\node[
    anchor=north,
    align=center,
    font=\fuentePortada\fontsize{10.0}{11.9}\selectfont
]
at ([xshift=12.10cm,yshift=-21.92cm]current page.north west)
{
    AYUDANTE DE LABORATORIO:
};

\node[
    anchor=north,
    align=center,
    text width=12.2cm,
    font=\fuentePortada\mdseries\fontsize{9.8}{11.6}\selectfont
]
at ([xshift=12.10cm,yshift=-22.44cm]current page.north west)
{
    \ayudanteLaboratorio
};


% ============================================================
% ESCUDO FACULTAD DE CIENCIAS
% ============================================================

\IfFileExists{imagenes/ciencias.png}{

    \node[
        anchor=north west,
        inner sep=0pt
    ]
    at ([xshift=3.49cm,yshift=-22.20cm]current page.north west)
    {
        \includegraphics[width=2.65cm]{imagenes/ciencias.png}
    };

    % ========================================================
    % LEYENDA DEBAJO DEL ESCUDO DE CIENCIAS
    %
    % Se deja en dos comandos separados para que sea fácil cambiar,
    % por ejemplo, a otra facultad.
    % ========================================================

    % Primera línea de la leyenda:
    % centrada con el escudo y muy cerca de la imagen.
    \node[
        anchor=north,
        align=center,
        font=\fuentePortada\fontsize{4.05}{4.60}\selectfont
    ]
    at ([xshift=4.815cm,yshift=-25.12cm]current page.north west)
    {
        \leyendaEscudoLineaUno
    };

    % Segunda línea:
    % separada de forma independiente para que UNAM quede
    % pegado a FACULTAD DE CIENCIAS como en la referencia.
    \node[
        anchor=north,
        align=center,
        font=\fuentePortada\fontsize{3.70}{4.20}\selectfont
    ]
    at ([xshift=4.815cm,yshift=-25.33cm]current page.north west)
    {
        \leyendaEscudoLineaDos
    };

}{}


% ============================================================
% SEMESTRE
% ============================================================

\node[
    anchor=north,
    align=center,
    font=\fuentePortada\mdseries\fontsize{11.0}{13.0}\selectfont
]
at ([xshift=12.10cm,yshift=-24.12cm]current page.north west)
{
    \semestreFinal
};


% ============================================================
% LUGAR Y FECHA
%
% Se deja editable al inicio del archivo con \fechaPortada.
% ============================================================

\node[
    anchor=north,
    align=center,
    font=\fuentePortada\mdseries\fontsize{10.4}{12.4}\selectfont
]
at ([xshift=12.10cm,yshift=-24.78cm]current page.north west)
{
    \fechaPortada
};


\end{tikzpicture}


\end{titlepage}


% ============================================================
% REINICIO DE NUMERACIÓN
% ============================================================

\setcounter{page}{2}

\hypersetup{pageanchor=true}

\fi


% ============================================================
% NOMBRES EN ESPAÑOL
% ============================================================

\renewcommand{\contentsname}{Índice}

\renewcommand{\tablename}{Tabla}

\renewcommand{\figurename}{Figura}

\renewcommand{\refname}{Bibliografía}


% ============================================================
% ÍNDICE
% ============================================================

\tableofcontents

\clearpage
\section{Implementación Original del Sistema}
\label{sec:org984960d}

Dado el modelo del sistema de estacionamiento, se plantearán las pruebas necesarias para verificar que el comportamiento de la implementación original del sistema sea el esperado por el modelo del mismo.

\vspace{\baselineskip}

En caso de que la implementación original del sistema no pase las pruebas planteadas, modificaremos el código para que este último tenga el comportamiento esperado por el modelo del sistema.

\vspace{\baselineskip}


\begin{lstlisting}[style=practica]
    ...

    TARIFA_BOLETO_PERDIDO = 300.0

    def calcular_total(self, minutos: int, tipo_cliente: str = "normal",
                       boleto_perdido: bool = False) -> float:
        """
        Calcula el total a pagar de acuerdo con el modelo de la práctica.

        Nota para el alumno:
	trate esta implementación como un sistema que debe verificarse.
	No asuma que todo lo que hace el código es correcto.
	"""
    
        if minutos < 0:
            raise ValueError("Los minutos no pueden ser negativos")


        if tipo_cliente not in ("normal", "frecuente"):
            raise ValueError("Tipo de cliente no válido")


        if boleto_perdido:
            total = self.TARIFA_BOLETO_PERDIDO
            

        elif minutos < 15:
            total = 0.0
            

        elif minutos <= 60:
            total = 20.0

	
        else:
            horas_adicionales = (minutos - 60) // 60
            total = 20.0 + (horas_adicionales * 15.0)

	
        if tipo_cliente == "frecuente":
            total *= 0.90

        return round(total, 2)

    ...
\end{lstlisting}
\section{Reglas del Sistema: Ticket de Estacionamiento}
\label{sec:org027dda5}

Las reglas del Sistema de Ticket de Estacionamiento son:

\begingroup
\setlength{\parskip}{0.85em}

\textbf{R1.} Los minutos deben ser un número entero mayor o igual que 0.

\textbf{R2.} Los tipos de clientes válidos son:

\begin{itemize}
\item \textbf{"normal"}
\item \textbf{"frecuente"}
\end{itemize}

\textbf{R3.} Si el boleto está perdido, el costo final es de \$300.00, sin importar los minutos ni el tipo de cliente.

\textbf{R4.} Si los minutos son 0, el costo es \$0.00

\textbf{R5.} De 1 a 15 minutos, el costo es \$0.00

\textbf{R6.} De 16 a 60 minutos, el costo es de \$20.00

\textbf{R7.} Después de 60 minutos se cobran \$20.00 por la primera hora y \$15.00 por cada hora adicional \textbf{iniciada}.

\textbf{R8.} Un cliente \textbf{"frecuente"} recibe un 10\% de descuento sobre el costo calculado, excepto cuando perdió el boleto.

\textbf{R9.} El costo nunca puede ser negativo.

\textbf{R10.} Si los minutos son negativos, el sistema debe lanzar un error.

\textbf{R11.} Si el tipo de cliente no es válido, el sistema debe lanzar un error.

\endgroup
\section{Construcción de la Tabla de casos}
\label{sec:org6eb7c2b}

\subsection{Planteamiento de pruebas}
\label{sec:orga7aacca}

Planteamos para el Sistema de Ticket de Estacionamiento:

\begin{itemize}
\item 3 casos normales;
\item 6 casos frontera;
\item 2 entradas inválidas;
\item 2 casos de interacción entre reglas;
\item 1 caso para números decimales en el resultado;
\item 1 caso adicional propuesto por el equipo (entrada inválida).
\end{itemize}
\subsection{Tabla de casos}
\label{sec:org87be2b3}


% ============================================================
% TABLA DE CASOS MULTIPÁGINA
%
% Esta tabla:
% - continúa automáticamente en la siguiente página;
% - repite el encabezado en las páginas siguientes;
% - deja espacio horizontal y vertical dentro de las celdas;
% - alinea solo Núm. al centro y el resto a la izquierda;
% - ajusta el texto normal dentro de cada columna;
% - nunca parte un caso entre dos páginas: cambia de página entre filas.
%
% IMPORTANTE:
% longtable puede dividir la TABLA entre páginas, pero una sola
% fila no debe contener una página completa de texto. Para los
% casos de prueba normales (una o varias líneas por celda)
% funciona correctamente.
% ============================================================

\begingroup

\small

% Espacio horizontal entre el texto y las líneas verticales.
\setlength{\tabcolsep}{8pt}

% Espacio vertical simétrico dentro de todas las celdas.
% makecell agrega la misma separación arriba y abajo del texto.
\renewcommand{\arraystretch}{1.10}
\setcellgapes{4pt}
\makegapedcells

% Centrar la tabla.
\setlength\LTleft{\fill}
\setlength\LTright{\fill}

\begin{longtable}{
    |C{1.00cm}
    |L{1.95cm}
    |L{2.30cm}
    |L{2.00cm}
    |L{1.55cm}
    |L{4.35cm}|
}

% ------------------------------------------------------------
% ENCABEZADO DE LA PRIMERA PÁGINA
% ------------------------------------------------------------

\hline

\textbf{Núm.} &
\textbf{Caso} &
\textbf{Entrada} &
\textbf{Resultado esperado} &
\textbf{Regla(s)} &
\textbf{Justificación}
\\

\hline

\endfirsthead


% ------------------------------------------------------------
% ENCABEZADO DE LAS PÁGINAS SIGUIENTES
% ------------------------------------------------------------

\hline

\textbf{Núm.} &
\textbf{Caso} &
\textbf{Entrada} &
\textbf{Resultado esperado} &
\textbf{Regla(s)} &
\textbf{Justificación}
\\

\hline

\endhead


% No agregamos texto extra al pie de cada fragmento.
\endfoot

% Tampoco agregamos una línea adicional al final:
% cada caso ya termina con \hline.
\endlastfoot


% ------------------------------------------------------------
% CASOS
% ------------------------------------------------------------


1 &
Normal &
10 minutos

"normal"

con boleto &
\$0.00 &
R5 &
Analiza el comportamiento del sistema entre 0 a 15 minutos.
\\
\hline

2 &
Normal &
40 minutos

"normal"

con boleto &
\$20.00 &
R6 &
Analiza el comportamiento del sistema entre 16 a 60 minutos.
\\
\hline

3 &
Normal &
1250 minutos

"normal"

con boleto &
\$320.00 &
R7 &
Analiza el comportamiento del sistema entre horas completas e inciadas.
\\
\hline

4 &
Frontera &
15 minutos

"normal"

con boleto &
\$0.00 &
R5 &
Analiza el comportamiento del sistema en el minuto 15, donde el costo permanece en \$0.00 antes de cambiar en el minuto 16 a \$20.00
\\
\hline

5 &
Frontera &
16 minutos

"normal"

con boleto &
\$20.00 &
R6 &
Analiza el comportamiento del sistema en el minuto 16, donde el costo cambia a \$20.00 después del minuto 15, donde estaba a \$0.00
\\
\hline

6 &
Frontera &
60 minutos

"normal"

con boleto &
\$20.00 &
R6 &
Analiza el comportamiento del sistema en el minuto 60, donde el costo permanece en \$20.00 antes de cambiar en el minuto 61 a \$35.00 por las horas iniciadas.
\\
\hline

7 &
Frontera &
61 minutos

"normal"

con boleto &
\$35.00 &
R7 &
Analiza el comportamiento del sistema en el minuto 61, donde el costo cambia a \$35.00 por las horas inciadas después del minuto 60, donde estaba a \$20.00
\\
\hline

8 &
Frontera &
120 minutos

"normal"

con boleto &
\$35.00 &
R7 &
Analiza el comportamiento del sistema en el minuto 120, donde el costo permanece en \$35.00 por las horas completas antes de cambiar en el minuto 121 a \$50.00 por las horas inciadas.
\\
\hline

9 &
Frontera &
121 minutos

"normal"

con boleto &
\$50.00 &
R7 &
Analiza el comportamiento del sistema en el minuto 121, donde el costo cambia a \$50.00 por las horas iniciadas después del minuto 120, donde estaba a \$35.00 por las horas completas.
\\
\hline

10 &
Entrada inválida &
-1 minuto

"normal"

con boleto &
Lanza

Excepción:

Minutos

inválidos. &
R1 y R10 &
Analiza el comportamiento del sistema cuando recibe numeros negativos, donde debería rechazarlos.
\\
\hline

11 &
Entrada inválida &
5 minutos

"premium"

con boleto &
Lanza

Excepción:

Tipo de

cliente

inválido. &
R2 y R11 &
Analiza el comportamiento del sistema cuando recibe tipos de clientes inexistentes, donde debería rechazarlos.
\\
\hline

12 &
Interacción entre reglas &
1250 minutos

"frecuente"

con boleto &
\$288.00 &
R7 y R8 &
Analiza el comportamiento del sistema cuando un cliente "frecuente" con boleto aplica su descuento del 10\% dados los minutos transcurridos.
\\
\hline

13 &
Interacción entre reglas &
1250 minutos

"frecuente"

sin boleto &
\$300.00 &
R3 y R8 &
Analiza el comportamiento del sistema cuando un cliente "frecuente" no tiene boleto, en cuyo no aplica su decuento del 10\% y se le cobra \$300.00.
\\
\hline

14 &
Números decimales en el resultado &
61 minutos

"frecuente"

con boleto &
\$31.50 &
R7 y R8 &
Analiza el comportamiento del sistema cuando se le cobra al cliente "frecuente" con boleto por 61 minutos transcurridos, en cuyo caso debe cobrar lo justo, que es \$31.50.
\\
\hline

15 &
Entrada inválida &
15.5 minutos

"normal"

con boleto &
Lanza

Excepción:

Minutos

inválidos. &
R1 &
Analiza el comportamiento del sistema cuando recibe minutos en decimales, donde debería rechazarlos.
\\
\hline


% ------------------------------------------------------------
% TÍTULO DE LA TABLA
% ------------------------------------------------------------

\caption{Casos establecidos para el Sistema Ticket de Estacionamiento.}
\label{tab:casos-ticket}

\end{longtable}

\nomakegapedcells

\endgroup
\section{Evidencia de Ejecución}
\label{sec:org9275fa5}

\subsection{Casos Normales: 1, 2 y 3}
\label{sec:orgd8fff6c}

\subsubsection{Implementación de la Prueba en Python}
\label{sec:org08d0099}

Prueba encargada de verificar el comportamiento del sistema en los casos normales 1, 2 y 3.

\vspace{\baselineskip}


\begin{lstlisting}[style=practica]
    # Arrange
    @pytest.mark.parametrize(
        "minutos, esperado",
        [
                            # | Caso | Descripción                                |
            (10, 0.00),     # |  1   | Cliente "normal" con boleto pasó 10 min.   |
            (40, 20.00),    # |  2   | Cliente "normal" con boleto pasó 40 min.   |
            (1250, 320.00), # |  3   | Cliente "normal" con boleto pasó 1250 min. |
        ],
    )
    def test_cliente_Normal_con_Boleto(sistema, minutos, esperado):
        """
        Verifica el cobro calculado para un cliente "normal" con boleto
        dado un cierto tiempo en minutos.
        """

        # Act
        resultado = sistema.calcular_total(minutos, "normal", False)

        # Assert
        assert resultado == esperado
\end{lstlisting}
\subsubsection{Resultado de la Prueba}
\label{sec:org53017db}

Los casos 1 y 2 pasaron la prueba; sin embargo el caso 3 la falló. Dado que el planteamiento de las pruebas esta basado en una interpretación correcta del modelo, analizamos el Código original.
\subsubsection{Código Original}
\label{sec:org3432ce7}


\begin{lstlisting}[style=practica]
    ...
    else:
        horas_adicionales = (minutos - 60) // 60
        total = 20.0 + (horas_adicionales * 15.0)
    ...
\end{lstlisting}

Al revisar el Código original, notemos que el cálculo realizado para conocer el cobro del estacionamiento con horas adicionales es incorrecto, ya que se pierde información de las horas iniciadas si solo se toma la parte entera de la división entre 60 minutos, pues si la divisón es exacta, tenemos horas completas, de lo contrario, ya inicio otra hora y se debe añadir una hora más al conteo de horas adicionales.
\subsubsection{Reparación del Código Original}
\label{sec:orgc780a7c}


\begin{lstlisting}[style=practica]
    ...
    else:
        horas_adicionales = math.ceil((minutos - 60) / 60) # Modificación
        total = 20.0 + (horas_adicionales * 15.0)
    ...
\end{lstlisting}

Para reparar el código correctamente conforme al modelo, debemos tomar el techo del resultado de la división para contemplar verdaderamente el número de horas adicionales.
\subsection{Casos Frontera: 4, 5, 6, 7, 8 y 9}
\label{sec:org7ac83df}

\subsubsection{Implementación de la Prueba en Python}
\label{sec:orgaa38c0f}

Prueba encargada de verificar el comportamiento del sistema en los casos frontera 4, 5, 6, 7, 8 y 9.

\vspace{\baselineskip}


\begin{lstlisting}[style=practica]
    # Arrange
    @pytest.mark.parametrize(
        "minutos, esperado",
        [
                          # | Caso | Descripción                              |
            (15, 0.00),   # |  4   | Último valor entre el rango de 1-15.     |
            (16, 20.00),  # |  5   | Primer valor entre el rango de 16-60.    |
            (60, 20.00),  # |  6   | Último valor entre el rango de 16-60.    |
            (61, 35.00),  # |  7   | Una hora completa y una hora iniciada.   |
            (120, 35.00), # |  8   | Dos horas completas sin hora iniciada.   |
            (121, 50.00), # |  9   | Dos horas completas y una hora iniciada. |
        ],
    )
    def test_cliente_Normal_con_Boleto_Casos_Frontera(sistema, minutos, esperado):
        """
        Verifica el cobro calculado para un cliente "normal" con boleto dado un cierto tiempo,
        en escencia cuando hay un cambio entre los límites de rangos de tiempo.
        """
    
        # Act
        resultado = sistema.calcular_total(minutos, "normal", False)
    
        # Assert
        assert resultado == esperado
\end{lstlisting}
\subsubsection{Resultado de la Prueba}
\label{sec:org8e1659b}

Los casos 5, 6, 7, 8 y 9 pasaron la prueba; sin embargo el caso 4 la falló. Dado que el planteamiento de las pruebas esta basado en una interpretación correcta del modelo, analizamos el Código original.
\subsubsection{Código Original}
\label{sec:orga8a4ccb}


\begin{lstlisting}[style=practica]
    ...
    elif minutos < 15:
        total = 0.0
    ...
\end{lstlisting}

Al revisar el Código original, notemos que el cálculo realizado para conocer el cobro del estacionamiento en 15 minutos es incorrecto, ya que no se toma en cuenta el minuto 15 para cobrar la cantidad correcta de \$0.00.
\subsubsection{Reparación del Código Original}
\label{sec:org5e2174e}


\begin{lstlisting}[style=practica]
    ...
    elif minutos <= 15: # Modificación
        total = 0.0
    ...
\end{lstlisting}

Para reparar el código correctamente conforme al modelo, debemos tomar en cuenta al minuto 15 dentro del rango correspondiente y cobrar \$0.00.
\subsection{Casos Entrada Inválida: 10, 11 y 15}
\label{sec:orgcffe898}

\subsubsection{Implementación de la Prueba en Python}
\label{sec:org2662b5c}

Prueba encargada de verificar el comportamiento del sistema en los casos de entrada inválida 10, 11 y 15, este último propuesto por el equipo.

\vspace{\baselineskip}

Antes de escribir la prueba, recordemos que tener excepciones muy generales para los "errores" que maneje el programa, indican una mala práctica de programación, ya que dificulta la documentación y posterior mantenimiento del programa.

\vspace{\baselineskip}

Por lo anterior, buscamos las excepciones correspondientes que debería lanzar el programa dados los casos  de entradas inválidas que hemos establecido, pero dado que no las encontramos, procedemos a definir excepciones personalizadas que documenten el caso correpondiente:

\begin{itemize}
\item \textbf{MinutosInvalidos:} Lanzada cuando no se cumple \textbf{R1} y \textbf{R10}.
\item \textbf{ClienteInvalido:} Lanzada cuando no se cumple \textbf{R2}.
\end{itemize}

Con ello, definimos las pruebas con las nuevas excepciones que esperamos ver.

\vspace{\baselineskip}


\begin{lstlisting}[style=practica]
    ...    
    from MinutosInvalidos import MinutosInvalidos
    from ClienteInvalido import ClienteInvalido
    ...

    
    # Test de Entradas inválidas.

    # Arrange
    @pytest.mark.parametrize(
        "minutos",
	[
                    # | Caso | Descripción                   |
            -1,     # | 10   | -1 min no están definidos.    |
            15.5,   # | 15   | 15.5 min no están definidos.  |
	],
    )
    def test_Minutos_Invalidos(sistema, minutos):
        """
	Verifica que el programa rechaze los Minutos Inválidos establecidos
	en la R1.
	"""
	# Siguiendo la R10, verifcamos que rechaze los minutos negativos.
	# Act and Assert
        with pytest.raises(MinutosInvalidos):
            sistema.calcular_total(minutos, "normal", False)

    def test_Cliente_Invalido(sistema):
        """
	Verifica que el programa rechaze los clientes que no están
	definidos en la R2. Siguiendo el procedimiento de R11.
	"""
	# | Caso | Descripción                                |
        # | 10   | Cliente "premium" no está bien definido.   |
	
	# Arrange, Act and Assert
	with pytest.raises(ClienteInvalido):
            sistema.calcular_total(5, "premium", False)
    ...
\end{lstlisting}
\subsubsection{Resultado de la Prueba}
\label{sec:orgca6e7ad}

Los casos 10, 11 y 15 fallaron la prueba; dado que el planteamiento de las pruebas esta basado en una interpretación correcta del modelo, analizamos el Código original.
\subsubsection{Código Original}
\label{sec:org0240896}


\begin{lstlisting}[style=practica]
    ...
    if minutos < 0:
        raise ValueError("Los minutos no pueden ser negativos")

    if tipo_cliente not in ("normal", "frecuente"):
       raise ValueError("Tipo de cliente no válido")
    ...
\end{lstlisting}

\vspace{\baselineskip}

Al revisar el Código original, notamos que para el caso 10 y 11 no se lanzan las excepciones que esperabamos porque la implementación original reporta las entradas inválidas con la excepción ValueError, pero como mencionamos anteriormente, esta es una mala práctica de programación. Por otro lado, en el caso 15, la implementación original no considera a los decimales como valor de entrada inválida para los minutos, cuando debería hacerlo por el modelo establecido.
\subsubsection{Reparación del Código Original}
\label{sec:org3556288}


\begin{lstlisting}[style=practica]
    ...
    if minutos < 0:
        raise MinutosInvalidos("Los minutos no pueden ser negativos") # Modificacion

    if isinstance(minutos, float): # Modificación
        raise MinutosInvalidos("Los minutos no pueden ser en decimal")
        
    if tipo_cliente not in ("normal", "frecuente"):
        raise ClienteInvalido("Tipo de cliente no válido") # Modificacion
    ...
\end{lstlisting}

Para reparar el código correctamente conforme al modelo, modificamos las excepciones que se lanzan cuando hay valores de entrada inválidos, para los minutos la excepción MinutosInvalidos y para los tipos de cliente ClienteInvalido, además, agregamos el caso donde consideramos a los decimales como valor inválido para los minutos.
\subsection{Casos Interacción entre reglas: 12 y 13}
\label{sec:orgc8b4d11}

\subsubsection{Implementación de la Prueba en Python}
\label{sec:orgfa065e8}

Prueba que verifica el comportamiento del sistema en los casos de interacción entre reglas 12 y 13.

\vspace{\baselineskip}


\begin{lstlisting}[style=practica]
    ...
    # Arrange
    @pytest.mark.parametrize(
        "boleto, esperado",
        [
                              # | Caso | Descripción                                  |
            (False, 288.00),  # | 12   | Cliente "frecuente" con boleto con 1250 min. |
            (True, 300.00),   # | 13   | Cliente "frecuente" sin boleto con 1250 min. |
        ],
    )

    def test_cliente_Frecuente_Casos_Interaccion_entre_reglas(sistema, boleto, esperado):
        """
        Verifica que el programa calcule correctamente el costo del estacionamiento
        dado que el cliente es frecuente, pase 1250 minutos y tenga o no boleto.
        """
        # Act
        resultado = sistema.calcular_total(1250, "frecuente", boleto)


	
        # Assert
        assert resultado == esperado
    ...
\end{lstlisting}
\subsubsection{Resultado de la Prueba}
\label{sec:orgd113ab8}

El caso 12 pasó la prueba, pero el 13 la falló; dado que el planteamiento de las pruebas esta basado en una interpretación correcta del modelo, analizamos el Código original.
\subsubsection{Código Original}
\label{sec:orge0e1e0f}


\begin{lstlisting}[style=practica]
    ...
    if boleto_perdido:
        total = self.TARIFA_BOLETO_PERDIDO     
    ...
    else:
        horas_adicionales = math.ceil((minutos - 60) / 60)
        total = 20.0 + (horas_adicionales * 15.0)

    if tipo_cliente == "frecuente":
        total *= 0.90
	
    return round(total, 2)
    ...
\end{lstlisting}

Al revisar el Código original, notamos que el flujo para asignarle el costo a un cliente que perdió su boleto está incorrecto, pues si bien asigna la tarifa de \$300.00, no termina el cálculo y prosigue a continuar evaluando casos, como el tipo de cliente, cuando en realidad la R3 prevalece sobre la R8.
\subsubsection{Reparación del Código Original}
\label{sec:org292c3f5}


\begin{lstlisting}[style=practica]
    ...
    if boleto_perdido:
        return self.TARIFA_BOLETO_PERDIDO # Modificación.        
    ...
    else:
        horas_adicionales = math.ceil((minutos - 60) / 60)
        total = 20.0 + (horas_adicionales * 15.0)
	
    if tipo_cliente == "frecuente":
        total *= 0.90

    return round(total, 2)
    ...
\end{lstlisting}

Para reparar el código correctamente conforme al modelo, modificamos el flujo; si el cliente perdió su boleto, simplemente retornamos el costo del estacionamiento de \$300.00 sin pasar por los otros casos.
\subsection{Caso Números decimales en el resultado: 14}
\label{sec:orgbce9df4}

\subsubsection{Implementación de la Prueba en Python}
\label{sec:org772b4d9}

Prueba que verifica el comportamiento del sistema en el caso 14 de números decimales en el resultado.

\vspace{\baselineskip}


\begin{lstlisting}[style=practica]
    ...
    # Arrange
    @pytest.mark.parametrize(
        "minutos, esperado",
        [
                          # | Caso | Descripción                                      |
            (61, 31.50),  # | 14   | Cliente "frecuente" con boleto paso 61 minutos.  |
        ],
    )
    def test_cliente_Frecuente_Casos_Interaccion_entre_reglas(sistema, minutos, esperado):
        """
        Verifica que el programa calcule correctamente el costo del estacionamiento
        dado que el cliente es frecuente y tenga boleto.
        """
        # Act
        resultado = sistema.calcular_total(minutos, "frecuente", False)
    
        # Assert
        assert resultado == pytest.approx(esperado)
	...
\end{lstlisting}

En este caso, podemos ver que el resultado proviene de operaciones con float, y estas pueden aparecer con pequeñas diferencias de precisión por el manejo de este tipo de dato en el lenguaje, por ello, ocupamos pytest.approx para comparar valores con cierta tolerancia que definimos como correctos. En los anteriores casos no es necesario porque obtenemos resultados enteros que si pueden compararse exactamente con el operador "==" sin utilizar pytest.approx.
\subsubsection{Resultado de la Prueba}
\label{sec:org610855a}

El caso 14 pasó la prueba; dado que el planteamiento de las pruebas esta basado en una interpretación correcta del modelo, determinamos que el Código original está bien planteado para este caso.
\subsubsection{Código Original}
\label{sec:orgf420f59}


\begin{lstlisting}[style=practica]
    ...
    def calcular_total(self, minutos: int, tipo_cliente: str = "normal",
                       boleto_perdido: bool = False) -> float:
        ...
        if boleto_perdido:
            return self.TARIFA_BOLETO_PERDIDO
	...
        else:
            horas_adicionales = math.ceil((minutos - 60) / 60)
            total = 20.0 + (horas_adicionales * 15.0)
	    
        if tipo_cliente == "frecuente":
            total *= 0.90

        return round(total, 2)
    ...
\end{lstlisting}

Al revisar el Código original y compararlo con el modelo del sistema, round funciona bien para determinar el costo de la tarifa en centavos. Pero en este caso, se vuelve un poco redundante al trabajar con tarifas que solo tengan como decimal a 0.0 o 0.5 y estos representarse de forma exacta en binario; sin embargo, en caso de que llegasen a cambiar los precios posteriormente, entonces sería de más utilidad dependiendo de como se establezca el redondeo de tarifas en el nuevo modelo.
\section{Preguntas de análisis}
\label{sec:org4cd9fc3}

\textbf{1.} ¿Qué diferencia hay entre un caso normal y un caso frontera?

R


\textbf{2.} ¿Cuál de sus pruebas considera más importante y por qué?

R


\textbf{3.} ¿Encontró algún comportamiento de la implementación que no coincida con el modelo?

R


\textbf{4.} ¿Una suite con 100 \% de pruebas aprobadas demuestra que el programa es correcto? Explique.

R


\textbf{5.} Si una IA generara automáticamente 50 pruebas, ¿qué tendría que revisar una persona antes de confiar en ellas?

R
\end{document}
